\documentclass[10pt,a4paper]{article}
\usepackage{amsmath,amssymb,amsthm,graphicx}
\usepackage[no-math]{fontspec}
\newenvironment{authorclaim}[1]{\par\medskip\noindent\textbf{#1.}\ }{\par\medskip}
\allowdisplaybreaks
\title{Randomness-efficient almost-universal hashing with near-2n circuit complexity}
\author{Lijie Chen, Jiatu Li, Tianqi Yang}
\date{doi:10.4230/LIPIcs.CCC.2022.23; CC BY 4.0}
\begin{document}
\maketitle
\noindent\textit{Editorial scope.} Counted once: the complete final hash construction Theorem3.9, incorporating all parameter/circuit regimes. Internal construction components and uniformity/explicitness continuation remain within the same item. Full original local proofs, exact imported background and complete AppendixA retained. Author responsibility for mathematical statements remains unchanged; no proof supplementation or author repair.\par\medskip
\noindent\textit{Allowed standard prerequisites (editorial).}\begin{itemize}
\item Almost-independent bits/permutations, biased sets, expander walks and imported high-girth/finite-field results. Exact published Theorems 2.4, 2.6, 2.8, 2.10, 2.11 and A.1, A.2 retained with full 43-entry bibliography. Source-local adaptation of high-girth graph explicitness fully included, not treated as an imported proof.
\end{itemize}\medskip
% Faithful typesetting transcription of the formal CCC2022.23 publication, Sections2 and AppendixA.
\section*{Author prerequisites (Section2)}
\paragraph{Notation.} We define $[n]\triangleq\{0,1,\ldots,n-1\}$, $B_{n,m}$ as the set of all functions from $\mathbb F_2^n$ to $\mathbb F_2^m$ and $B_n\triangleq B_{n,1}$. The Hamming weight of $x\in\mathbb F_2^n$, denoted by $|x|$, is defined as the number of 1's in $x$; the Hamming distance $\Delta(x,y)$ of $x$ and $y$ from $\mathbb F_2^n$ is defined as the Hamming weight of the bitwise XOR of them. The concatenation of two strings $x$ and $y$ is denoted by $x\Vert y$, and the $i$th bit of $x$ (0-indexed) is denoted by $x_i$. Graphs are undirected by default. We assume that all functions used as parameters of our constructions are poly-logarithmic time constructible, i.e., for any function $\ell(n)$ that is used as a parameter in our results, there is a polynomial-time algorithm $A$ such that $A(n)$ outputs the binary representation of $\ell(n)$.
We use $x\leftarrow D$ to denote a random element $x$ sampled from a distribution $D$, and $D(x_0)\triangleq\Pr_{x\leftarrow D}[x=x_0]$. Natural numbers are represented in binary when being fed into Turing machines or circuits; and $1^n$ represents the unary representation of $n$. We assume basic familiarity with cryptographic primitives like one-way functions and pseudorandom functions, and typical complexity classes like $\mathsf P$, $\mathsf{NP}$, and $\oplus\mathsf P$ (see, e.g., [3]).
\subsection*{2.1 Probability Theory}
Let $U(S)$ be the uniform distribution supported on a set $S$, and $U_\ell\triangleq U(\{0,1\}^\ell)$. A family $D=\{D_n\}_{n\ge1}$ of distributions is said to be samplable with seed length $\ell(n)$ (or $\ell(n)$-samplable) if there is a polynomial-time algorithm $A$ such that $A(1^n,U_{\ell(n)})$ samples $D_n$. For every $s\in\{0,1\}^\ell$, we say $A(1^n,s)$ is the element corresponding to the seed $s$.
The statistical distance between two distributions $D_1$ and $D_2$ over a set $S$, denoted by $\operatorname{SD}(D_1,D_2)$, is defined as
\[\operatorname{SD}(D_1,D_2)\triangleq\frac12\sum_{u\in S}|D_1(u)-D_2(u)|.\]
A distribution $D_2$ is said to be $\delta$-close to $D_1$ if their statistical distance is at most $\delta$. Statistical distance characterizes the intractability of distinguishing two distributions by any test (even computationally unbounded): if $D_1$ and $D_2$ have statistical distance at most $\delta$, then for any stochastic process $T$, $|\Pr_{u\leftarrow D_1,T}[T(u)=1]-\Pr_{u\leftarrow D_2,T}[T(u)=1]|\le\delta$.
\subsection*{2.2 Circuit Classes}
In this paper we will work with various circuit classes. In general, a circuit is an acyclic graph where each of its nodes can be an input variable, a constant $c\in\{0,1\}$, or a gate. One or more nodes are marked as output nodes (together with an index $i$ denoting the corresponding output bit). The depth of a circuit is the number of edges on the longest path from any input variable to an output node. An $n$-input $m$-output circuit computes a function in $B_{n,m}$.
\begin{description}
\item[{$B_2$ circuits.}] A $B_2$ circuit (or general circuit) contains fan-in-2 gates that can compute any binary function $f\in B_2$. The size of a $B_2$ circuit refers to the number of gates involved.
\item[{$\mathsf{NC}^1$ circuits.}] An $\mathsf{NC}^1$ circuit is a $B_2$ circuit with $O(\log n)$ depth. Note that a single-output $\mathsf{NC}^1$ circuit of depth $d$ can be converted into a formula of size $O(2^d)$.
\item[{$\mathsf{CC}^0[2]$ circuits.}] A $\mathsf{CC}^0[2]$ circuit is a constant-depth circuit with only XOR gates of unbounded fan-in. It is easy to see that $\mathsf{CC}^0[2]$ circuits can only compute linear functions over $\mathbb F_2$. The complexity of a $\mathsf{CC}^0[2]$ circuit is usually measured by the number of wires.
\item[{$\mathsf{AC}^0[m]$ circuits.}] An $\mathsf{AC}^0[m]$ circuit is a constant-depth circuit with fan-in-1 NOT gates and unbounded fan-in gates over $\{\mathrm{AND},\mathrm{OR},\mathrm{MOD}_m\}$, where $\mathrm{MOD}_m(x_1,\ldots,x_k)=1$ if and only if $x_1+x_2+\cdots+x_k\equiv0\pmod m$. Similar to $\mathsf{CC}^0[2]$ circuits, the complexity of an $\mathsf{AC}^0[m]$ circuit is measured by the number of wires.
\end{description}
A probabilistic $B_2$ circuit (or simply a probabilistic circuit) with input size $n$ and output size $m$ is a distribution $C_n$ over $n$-input $m$-output $B_2$ circuits. The circuit complexity of a probabilistic circuit is defined as the maximum complexity of functions in its support. A family $C=\{C_n\}_{n\ge1}$ of $n$-input 1-output probabilistic circuit is said to decide a language $L$ with error probability $\varepsilon=\varepsilon(n)$ if for sufficiently large $n$ and all $x\in\mathbb F_2^n$, $\Pr_{C\leftarrow C_n}[C(x)\ne L(x)]\le\varepsilon$.
\subsection*{2.3 Hash and 1-detector}
\begin{authorclaim}{Definition2.1 (Almost universal hash function)}
Let $m=m(n)$ and $\varepsilon=\varepsilon(n)$ be two parameters. An $m$-output $\varepsilon$-almost universal hash function is a family of distributions $H=\{H_n\}_{n\ge1}$, where $H_n$ is supported on $B_{n,m}$, such that $\Pr_{h\leftarrow H_n}[h(x)=h(y)]\le\varepsilon(n)$ for all $x\ne y$ and sufficiently large $n$. The parameter $\varepsilon$ is called its collision probability. It is linear if every function in the support of $H_n$ is linear. It is said to be $s(n)$-samplable if the family $H$ of distributions is samplable with seed length $s(n)$.
\end{authorclaim}
We will make heavy use of the notion of 1-detectors in [13], as parts of our almost universal hash construction.
\begin{authorclaim}{Definition2.2 (1-detector)}
An $m$-output (randomized) 1-detector with range $r$ and error $\varepsilon$ is a family of distributions $D=\{D_n\}_{n\ge1}$, where $D_n$ is supported on linear functions in $B_{n,m}$, such that for sufficiently large $n$, $\Pr_{L\leftarrow D_n}[L(x)\ne0]\le\varepsilon(n)$ for all $x\in\mathbb F_2^n$ satisfying $0<|x|\le r$. It is said to be $s(n)$-samplable if the family $D$ of distributions is samplable with seed length $s(n)$.
\end{authorclaim}
We emphasize that the definition requires the functions in the family of distributions to be linear. This ensure that, for any randomized 1-detector $D_n$ with range $r$ and error $\varepsilon$, and any $x_1\ne x_2$ with Hamming distance at most $r$, it holds that $\Pr_{L\leftarrow D_n}[L(x_1)\ne L(x_2)]\le\varepsilon(n)$. The shrinkage of a hash function (or 1-detector) is defined as the input length $n$ divided by the output length $m$ (e.g., poly-logarithmic shrinkage means $m=n/\operatorname{polylog}(n)$).
\subsection*{2.4 $\varepsilon$-biased set and $k$-wise independence}
To derandomize the low-complexity hash function in [13], we need several standard derandomization tools including strongly explicit $\varepsilon$-biased sets, $k$-wise independent distributions, and $k$-wise independent permutations.
\begin{authorclaim}{Definition2.3 ($\varepsilon$-biased set)}
For any $n\ge1$ and $\varepsilon\in(0,1/2)$, a set $S\subseteq\mathbb F_2^n$ is said to be $\varepsilon$-biased if for all non-zero $v\in\mathbb F_2^n$, $\Pr_{w\leftarrow S}[\langle w,v\rangle=0]\in[1/2-\varepsilon,1/2+\varepsilon]$.
\end{authorclaim}
\begin{authorclaim}{Theorem2.4 (Alon, Goldreich, H\aa stad, and Peralta [1])}
For any constant $\varepsilon\in(0,1/2)$, there is a family $\{S_n\subseteq\mathbb F_2^n\}_{n\ge1}$ of $\varepsilon$-biased sets such that $|S_n|=\widetilde O(n^2)$. Moreover, there is an algorithm $A(n,i,j)$ running in time $\operatorname{poly}(\log n)$ that computes the $j$th bit of the $i$th vector in $S_n$.
\end{authorclaim}
\begin{authorclaim}{Definition2.5}
A $k$-wise $\varepsilon$-dependent distribution with length $n$ and alphabet size $p$ is a distribution $D$ over $[p]^n$ such that for all $0\le i_1<i_2<\cdots<i_k<n$, the distribution over $[p]^k$ obtained by restricting $D$ to the $i_1$th, $i_2$th, $\ldots$, $i_k$th coordinates is $\varepsilon$-close to the uniformly distribution over $[p]^k$. It is said to be a $k$-wise independent distribution if $\varepsilon=0$.
\end{authorclaim}
\begin{authorclaim}{Theorem2.6 (Alon, Goldreich, H\aa stad, and Peralta [1])}
Let $\varepsilon=\varepsilon(n)>0$ be a parameter. There is an algorithm $A$ such that $A(1^n,k,U_r)$ runs in time $\operatorname{poly}(n,\log(1/\varepsilon))$ and samples a $k$-wise $\varepsilon$-dependent distribution with length $n$ and alphabet size 2 for every $1\le k\le n$, where $r=O(k+\log\log n+\log(1/\varepsilon))$. Moreover, there is an algorithm $B$ such that $B(n,k,v,i)$ runs in $\operatorname{poly}(\log n,k,\log(1/\varepsilon))$ time and computes the $i$th coordinate of the string corresponding to the seed $v$.
\end{authorclaim}
\begin{authorclaim}{Definition2.7}
Let $n\ge1$ be the number of elements, and $P_n$ be the set of all permutations over $[n]$. A $k$-wise $\varepsilon$-dependent permutation is a distribution $D$ over $P_n$ such that for all $0\le i_1<i_2<\cdots<i_k<n$, the distribution of $(f(i_1),f(i_2),\ldots,f(i_k))$ with $f\leftarrow D$ is $\varepsilon$-close to the uniform distribution over $\{(j_1,j_2,\ldots,j_k)\mid j_1,\ldots,j_k\in[n]\text{ are pairwisely distinct}\}$.
\end{authorclaim}
\begin{authorclaim}{Theorem2.8 (Kaplan, Naor, and Reingold [25], Theorem5.9)}
Let $n$ be a power of 2, $2\le k\le n$, and $\varepsilon>2^{-n}$. There is an $n$-element $k$-wise $\varepsilon$-dependent permutation that is samplable with seed length $O(k\log n+\log(1/\varepsilon))$. Moreover, there is an algorithm $A(n,v,i)$ that runs in $\operatorname{poly}(\log n,k,\log(1/\varepsilon))$ time and outputs $\rho_v(i)$ for the permutation $\rho_v$ corresponding to the seed $v$.
\end{authorclaim}
\subsection*{2.5 Expander Graphs}
We will need strongly explicit expander graphs in our proofs. We first recall the definition of expander graphs.
\begin{authorclaim}{Definition2.9}
An $n$-vertex $d$-regular graph $G$ is called an $(n,d,\lambda)$ expander graph (or $(n,d,\lambda)$-graph) if $\lambda_2(G)\le\lambda$, where $\lambda_2(G)$ denotes the second largest eigenvalue of normalized adjacency matrix (i.e., adjacency matrix divided by $d$) of $G$. A family of graphs $\{G_n\}_{n\ge1}$ is an expander graph family if there exists constants $d$ and $\lambda<1$ such that for all sufficiently large $n$, $G_n$ is an $(n,d,\lambda)$-graph.
\end{authorclaim}
We will make use of the following construction of strongly explicit expander graphs.
\begin{authorclaim}{Theorem2.10 (Strongly explicit expander; see [3, Theorem21.19])}
There exists an expander graph family $\{G_n\}_{n\ge1}$ and an algorithm $A(n,v,i)$ that runs in $\operatorname{polylog}(n)$ time (i.e., polynomial in input length) and outputs the $i$th neighbor of the node $v$ in $G_n$, where $v\in[n]$ and $i\in[d]$.
\end{authorclaim}
Performing random walk on expanders is a standard derandomization technique. Consider the task to find a good element from $n$ elements in which there is a constant fraction of them being good. A trivial approach is to sample $\ell$ independently random elements, which has $\exp(-\Omega(\ell))$ error probability but requires $\ell\log n$ random bits. By applying the following lemma, we can reduce the randomness complexity to $O(\log n+\ell)$ while keeping the error probability to be exponentially small.
\begin{authorclaim}{Lemma2.11 (Expander walk; see [3, Theorem21.12])}
Let $G$ be an $(n,d,\lambda)$-graph and $S$ be a subset of vertices of size at most $\beta n$ for some $\beta\in(0,1)$. Let $X_1,X_2,\ldots,X_k$ be random variables denoting a random walk in $G$ (where $X_1$ is uniformly chosen), then
\[\Pr[\forall1\le i\le k, X_i\in S]\le((1-\lambda)\sqrt\beta+\lambda)^{k-1}.\]
\end{authorclaim}
\subsection*{2.6 Graph with large girth}
The girth of an undirected graph is the length of the minimum cycle in it. We need the following construction of strongly explicit graphs with large girth.
\begin{authorclaim}{Theorem2.12 (Adapted from [27]; see AppendixA)}
Let $r=r(n)=n^{o(1)}$ be a parameter. For every sufficiently large $n$, there exists an $m=\Theta(n/r)$ and a regular graph $G_{m,n}$ with $m$ vertices, $n$ edges, and girth $\Omega(\log n/\log r)$. Moreover, there exists a $\operatorname{polylog}(n)$-time algorithm $A(n,i)$ for $i\in[n]$ that outputs the indices of the two endpoints of the $i$th edge in $G_{m,n}$, and a $\operatorname{polylog}(n)$-time algorithm $B(n,i,j)$ for $i\in[m]$ that outputs the $j$th edge attaching to the $i$th vertex.
\end{authorclaim}

% Faithful formal-publication transcription: Section3, pp23:15--23:24. Historical intuition is not needed as substitute for the full formal proofs below.
\section*{3 Randomness-efficient low-complexity hash functions}
In this section, we present constructions of randomness-efficient low-complexity hash functions with various parameters and properties.
In Section3.1, we show almost universal hash functions can be constructed from 1-detectors (Lemma3.2). In Section3.2 we give a derandomized version of construction of 1-detectors from [13] based on high-girth graphs (Lemma3.3), and use that to construct a low-complexity hash function with poly-logarithmic seed-length and $n/\operatorname{polylog}(n)$ output length (Theorem3.4).
In Section3.3, we show how to reduce the output length from $n/\operatorname{polylog}(n)$ to $\operatorname{polylog}(n)$, by composing almost universal hash families. Finally, in Section3.4 and3.5, we establish the uniformity and explicitness of our constructions, which are essential for our applications to hardness magnification and PRF constructions.
\paragraph{Notation.} In this section, for a construction $H$ with parameters $a,b,c$, we will use $H_{\hat a,\hat b,\hat c}$ to denote the specific construction $H$ with the parameters specified to $\hat a,\hat b$ and $\hat c$. When the parameters are obvious from the context, we often omit the subscripts and simply write it as $H$.
\subsection*{3.1 General construction from 1-detectors}
We first give a general construction of almost universal hash functions from 1-detectors, using the following sampling procedure.
\begin{authorclaim}{Lemma3.1}
For all integers $n,b$ and $r$ such that $b\le n$ and $\max\{10n/b,10\log\log n\}\le r\le n$, there is a distribution $D^{\mathrm{samp}}_{n,b,r}$ supported on $[n]^b$ samplable by $O(r\log(n/b))$ bits, such that the following conditions hold:
\begin{enumerate}
\item (Hitting Condition). For all $x\in\mathbb F_2^n$ with $|x|\ge r$, it holds that
\begin{equation}\tag{1}
\Pr_{(w_0,\ldots,w_{b-1})\leftarrow D^{\mathrm{samp}}_{n,b,r}}\left[\bigwedge_{j\in[b]}[x_{w_j}=0]\right]\le2\exp\left[-\frac{br}{2n}\right].
\end{equation}
\item (Ordering) For every $(w_0,\ldots,w_{b-1})\in\operatorname{supp}(D^{\mathrm{samp}}_{n,b,r})$, it holds that $w_0<w_1<\cdots<w_{b-1}$.
\item (Explicitness) There are algorithms $A^{\mathrm{samp}}_{n,b,r}(v,i)$ and $B^{\mathrm{samp}}_{n,b,r}(v,j)$ running in $\operatorname{poly}(\log n,r)$ time such that
\[A^{\mathrm{samp}}_{n,b,r}(v,i)\text{ outputs }\begin{cases}j&\text{if }w_j=i\text{ for some }j\in[b],\\\bot&\text{otherwise;}\end{cases}\]
\[B^{\mathrm{samp}}_{n,b,r}(v,j)\text{ outputs }w_j,\]
where $(w_0,\ldots,w_{b-1})\in\operatorname{supp}(D^{\mathrm{samp}}_{n,b,r})$ is the vector corresponding to the seed $v$.
\end{enumerate}
\end{authorclaim}
\begin{proof}
We firstly describe a sampling procedure that satisfies Equation(1) but has a long seed length and then reduce the seed length to $O(r\log(n/b))$ bits using the explicit $k$-wise $\varepsilon$-dependent distribution from Theorem2.6. The sampling procedure is simple: we first partition $[n]$ into $b$ consecutive groups $g_0,g_1,\ldots,g_{b-1}$ of size either $\lfloor n/b\rfloor$ or $\lceil n/b\rceil$; we then uniformly choose an $i_j\in g_j$ for each $j\in[b]$; finally we output the tuple $(i_0,i_1,\ldots,i_{b-1})$.
For any $x\in\mathbb F_2^n$ with Hamming weight at least $r$, for some $\ell\le r$, there are $\ell$ groups such that there are at least $r$ 1-indices (i.e., the corresponding bits of $x$ are 1) in these groups. Assume that there are $a_j$ 1-indices in the $j$th among these $\ell$ groups. The probability that none of the 1-indices is sampled is at most
\begin{equation}\tag{2}
\begin{aligned}
\prod_{j=1}^\ell\left[1-\frac{a_jb}{2n}\right]
&\le\exp\left(\sum_{j=1}^\ell\log\left[1-\frac{a_jb}{2n}\right]\right)\\
&\le\exp\left(-\sum_{j=1}^\ell\frac{a_jb}{2n}\right)\le\exp\left[-\frac{br}{2n}\right].
\end{aligned}
\end{equation}
A trivial implementation of the sampling procedure above needs seed length $O(b\log(n/b))$: for each group, we sample a random binary string of length $t\triangleq\lceil\log(n/b)\rceil$ indicating the index we want to choose. Therefore, we would need to sample $bt$ bits in total. The observation here is that the argument in (2) only involves $\ell\le r$ groups. Hence, the analysis still works if we sample those required $bt$ bits from an $rt$-wise $\exp(-br/(2n))$-dependent distribution with alphabet 2. Note that this incurs an additional error of $\exp(-br/(2n))$.
By Theorem2.6, the distribution is samplable with seed length
\[O(rt+\log\log(bt)+br/(2n))=O(r\log(n/b)),\]
and the algorithms $A^{\mathrm{samp}}_{n,b,r}$ and $B^{\mathrm{samp}}_{n,b,r}$ can be constructed straightforwardly from the algorithm $B$ in Theorem2.6.\footnote{To compute $A^{\mathrm{samp}}_{n,b,r}(v,i)$ for some $i$ belonging to the group $g_a$, we only need to check if $i$ is the selected element of the group $g_a$, and then output $a$ if the answer is yes, and output $\bot$ otherwise.}
\end{proof}
Now we are ready to give the construction of our hash function from 1-detectors. Recall that their definitions are given in Section2.3.
\begin{authorclaim}{Construction1 (Hash function $H^L$ from 1-detectors $L$)}
Let $n$ be the input length and $b(n)=o(n)$ be a parameter. Given a randomized 1-detector $L=\{L_n\}_{n\ge1}$ with range $\hat r$, we construct a hash function $H_b^L=\{H_n\}_{n\ge1}$ as follows.
\begin{itemize}
\item Let $D^{\mathrm{samp}}$ be the distribution in Lemma3.1 with parameters $n,b$ and $r=\hat r$. For each $L$ in the support of $L$ and each $D=(j_0,\ldots,j_{b-1})$ in the support of $D^{\mathrm{samp}}$, we define a function $h_{L,D}$ as
\[h_{L,D}(x)\triangleq L(x)\Vert x_{j_0}\Vert x_{j_1}\Vert\cdots\Vert x_{j_{b-1}}.\]
\item $H_n$ is then defined to be the distribution generated as follows: sample $L\leftarrow L_n$ and $D\leftarrow D^{\mathrm{samp}}$, and then output $h_{L,D}$.
\end{itemize}
\end{authorclaim}
We show that Construction1 indeed gives almost universal hash functions.
\begin{authorclaim}{Lemma3.2 ($H^L$ is a linear almost universal hash)}
Let $L=\{L_n\}_{n\ge1}$ be an $\hat m$-output $\hat s$-samplable randomized 1-detector with error $\hat\varepsilon$ and range $\hat r$, such that $\hat m=o(n)$ and $\hat r=\omega(\log n)$. Let $b=b(n)=o(n)$ be a parameter and $r=r(n)=\hat r(n)$. The followings hold for $H_b^L$:
\begin{enumerate}
\item $H_b^L$ has output length $m=\hat m+b$.
\item $H_b^L$ is $(\hat s+\hat r\log(n/b))$-samplable.
\item $H_b^L$ is a linear $\varepsilon$-almost universal hash function, where $\varepsilon\triangleq\hat\varepsilon+\exp(-\Omega(br/n))$.
\end{enumerate}
\end{authorclaim}
\begin{proof}
The first two items follow directly from the definition of $H^L$ from Construction1. So we will only establish the last item. We need to show that for any two inputs $x_1\ne x_2$ from $\{0,1\}^n$, their hash values collide with probability at most $\varepsilon$.
Consider the following two cases depending on whether the Hamming distance between $x_1$ and $x_2$ is small or large:
\begin{itemize}
\item If $0<\Delta(x_1,x_2)\le r$, then by the definition of randomized 1-detector (Definition2.2), $L(x_1)$ equals to $L(x_2)$ with probability at most $\hat\varepsilon$ for $L\leftarrow L_n$ (note that $L$ is a linear function over $\mathbb F_2$). Since the hash values contain the output of the 1-detector as the first $\hat m$ bits, the hash values collide with probability at most $\hat\varepsilon$.
\item If $\Delta(x_1,x_2)>r$, then $y\triangleq x_1\oplus x_2$ has Hamming weight at least $r$. By Lemma3.1, a random $D\leftarrow D^{\mathrm{samp}}_{n,b,r}$ fails to contain all 1-indices of $y$ with probability at most $\exp(-\Omega(br/n))$, which means that the last $b$ bits of the hash values of $x_1$ and $x_2$ collide with probability at most $\exp(-\Omega(br/n))$.
\end{itemize}
It follows immediately that the collision probability is at most $\varepsilon=\hat\varepsilon+\exp(-\Omega(br/n))$.
\end{proof}
\subsection*{3.2 Randomness-efficient low-complexity 1-detectors}
Now we will present the construction of randomness-efficient low-complexity 1-detectors. Applying Lemma3.2, this construction yields a hash function with inverse-super-polynomial collision probability, short seed length, and moderately large shrinkage.
\begin{authorclaim}{Construction2 (Construction of the 1-detector $L^{\mathrm{hg}}$ from high-girth graphs)}
Let $k=k(n)=n^{o(1)}$ be the shrinkage parameter, $\log n\le r(n)\le n/2$ be the range. Let $\gamma=\gamma(n)\triangleq\log^2n/\log k$. $L^{\mathrm{hg}}_{k,r}=\{L_n\}_{n\ge1}$ is constructed as follows.
\begin{itemize}
\item Let $n'$ be the smallest power of two that is larger than $n$. Let $G_{m,n'}$ be the regular graph with $m=\Theta(n'/k)$ vertices, $n'$ edges, and girth $g=\Theta(\log n/\log k)$ from Theorem2.12. Let $D=\{D_{n'}\}$ be the explicit family of $r$-wise $2^{-\gamma}$-dependent permutation in Theorem2.8. Let $\Gamma(i)$ be the set of indices of edges incident to the $i$th vertex in $G_{m,n'}$. For each permutation $\sigma\in\operatorname{supp}(D_{n'})$, we define a function $L_\sigma:\{0,1\}^n\to\{0,1\}^m$ as
\[L_\sigma(x)_i\triangleq\bigoplus_{j\in[n]\text{ s.t. }\sigma(j)\in\Gamma(i)}x_j\qquad\forall i\in[m].\]
\item $L_n$ is then defined to be the distribution generated as follows: sample $\sigma\leftarrow D_{n'}$ and then output $L_\sigma$.
\end{itemize}
\end{authorclaim}
\begin{authorclaim}{Lemma3.3 ($L^{\mathrm{hg}}$ is a 1-detector)}
Let $k=k(n)=n^{o(1)}$ be the shrinkage parameter, $\log n\le r(n)\le n/2$ be the range. The followings hold for $L^{\mathrm{hg}}_{k,r}=\{L_n\}_{n\ge1}$:
\begin{enumerate}
\item $L^{\mathrm{hg}}_{k,r}$ has output size $\Theta(n/k)$ and seed length $O(r\log n)$.
\item Every $L\in L_n$ can be computed by depth-1 $\mathsf{CC}^0[2]$ circuits of wire complexity at most $2n$, or $B_2$ circuits of size $2n$ and depth $\log k+O(1)$.
\item $L^{\mathrm{hg}}_{k,r}$ is a randomized 1-detector with range $r$ and error $\exp(-\Omega(\log^2n/\log k))$.
\end{enumerate}
\end{authorclaim}
\begin{proof}
Recall that $\gamma=\log^2n/\log k$ from Construction2. For the first item, the output size follows directly from the definition of $L^{\mathrm{hg}}_{k,r}$. The seed length of $L^{\mathrm{hg}}_{k,r}$ equals the seed length of the required $r$-wise $2^{-\gamma}$-dependent permutation, which is $O(r\log n+\gamma)=O(r\log n)$ from Theorem2.8 (note that $r\ge\log n$). The second item follows directly from the definition of $L_\sigma$ in Construction2 together with the fact that $G_{m,n'}$ has maximum degree $O(k)$ (since it is regular).
In the rest of the proof we establish the third item. We need to show that for any non-zero $x\in\mathbb F_2^n$, the probability that $L(x)=0$ with $L\leftarrow L_n$ is at most $\exp(-\Omega(\gamma))$. For any input $x$ with Hamming weight $0<|x|<g$ (where $g$ is the girth of $G_{m,n}$), we must have $L_\sigma(x)\ne0$ for all $\sigma\in D_{n'}$, since otherwise we can extract from the edges $\{\sigma(i)\mid x_i=1\wedge i\in[n]\}$ a cycle of size length less than $g$. Now we consider the case when $g\le|x|\le r$. Let $i_1,i_2,\ldots,i_{|x|}$ be the indices corresponding to the 1's in $x$ (i.e., $x_{i_j}=1$ for all $j$). Since our construction is linear, only these bits can influence the output. Also, by the $r$-wise $2^{-\gamma}$-dependence of $D_{n'}$, we have
\begin{align*}
\Pr_{\sigma\leftarrow D_{n'}}[L_\sigma(x)=0]
&=\Pr_{\sigma\leftarrow D_{n'}}\left[\bigwedge_{i'\in[m]}[|\{i_j\mid\sigma(i_j)\in\Gamma(i')\}|\text{ is even}]\right]\\
&\le\Pr_{S\subseteq[n'],\,|S|=|x|}\left[\bigwedge_{i'\in[m]}[|S\cap\Gamma(i')|\text{ is even}]\right]+2^{-\gamma}\\
&=\mathbb E_{\substack{S\subseteq[n']\\|S|=|x|-\alpha}}\left[\Pr_{\substack{S'\subseteq[n']\setminus S\\|S'|=\alpha}}\left[\bigwedge_{i'\in[m]}[|(S\cup S')\cap\Gamma(i')|\text{ is even}]\right]\right]+2^{-\gamma},
\end{align*}
where $\alpha\triangleq\lceil g/3\rceil$.
The final observation is that for any fixed $S$, there should be at most one $S'$ satisfying the condition in the summation. Indeed, if two sets $S'_1$ and $S'_2$ satisfy the condition at the same time, then the symmetric difference of them contains a cycle in $G_{m,n'}$ of length $2\alpha<g$ for sufficiently large $n$, contracting the fact that $G_{m,n'}$ has girth $g$. In particular, it means for every fixed $S$, we have
\[\Pr_{\substack{S'\subseteq[n']\setminus S\\|S'|=\alpha}}\left[\bigwedge_{i'\in[m]}[|(S\cup S')\cap\Gamma(i')|\text{ is even}]\right]\le\frac1{\binom{n'-|x|+\alpha}{\alpha}}.\]
Since $\alpha=\lceil g/3\rceil=\Theta(\log n/\log k)$, putting everything together, for sufficiently large $n$ we have
\[\Pr_{\sigma\leftarrow D_n}[L_\sigma(x)=0]\le\frac1{\binom{n'-|x|+\alpha}{\alpha}}+2^{-\gamma}=\exp(-\Omega(\gamma)).\]
\end{proof}
Plugging $L^{\mathrm{hg}}$ into Construction1 with appropriately chosen parameters, we immediately obtain a linear hash $H^{\mathrm{md}}$ with poly-logarithmic seed length and shrinkage as follows. (Here md stands for moderate, meaning that the hash has moderate (poly-logarithmic) shrinkage.)
\begin{authorclaim}{Construction3 (Construction of hash $H^{\mathrm{md}}$ with poly-logarithmic shrinkage)}
Let $\beta\in(0,2]$ be a constant and $\omega(\log n)\le\delta(n)\le O(\log^2n/\log\log n)$ be the error parameter. $H^{\mathrm{md}}_{\beta,\delta}=\{H_n\}_{n\ge1}$ is constructed as follows.
\begin{itemize}
\item Setting $k=\log^\beta n$ and $r=\delta\log^\beta n$, $L^{\mathrm{hg}}=L^{\mathrm{hg}}_{k,r}$ (from Construction2) is an $O(\delta\log^{\beta+1}n)$-samplable $\Theta(n/\log^\beta n)$-output randomized 1-detector with range $r$ and error $\exp(-\Omega(\delta))$.
\item Setting $b=n/\log^\beta n$, $H^{\mathrm{md}}$ is now defined to be $H_b^{L^{\mathrm{hg}}}$ (from Construction1).
\end{itemize}
\end{authorclaim}
The following theorem follows directly from the Construction3, Lemma3.2, and Lemma3.3.
\begin{authorclaim}{Theorem3.4 (Properties of the intermediate hash construction $H^{\mathrm{md}}$)}
Let $\beta\in(0,2]$ be the shrinkage parameter and $\omega(\log n)\le\delta(n)\le O(\log^2n/\log\log n)$ be the error parameter. The followings hold for $H^{\mathrm{md}}_{\beta,\delta}=\{H_n\}_{n\ge1}$:
\begin{enumerate}
\item $H^{\mathrm{md}}_{\beta,\delta}$ is a linear $O(\delta\log^{\beta+1}n)$-samplable $\Theta(n/\log^\beta n)$-output $\exp(-\Omega(\delta))$-almost universal hash function.
\item Every $H\in H_n$ can be computed either by a depth-1 $\mathsf{CC}^0[2]$ circuit of wire complexity $2n$, or by a $B_2$ circuit of size $2n$ and depth $\beta\log\log n+O(1)$.
\end{enumerate}
\end{authorclaim}

\subsection*{3.3 Shrinkage reduction of hash function}
Now we reduce the output length of the hash function $H^{\mathrm{md}}$ in Construction3 from $n/\operatorname{polylog}(n)$ to $\operatorname{polylog}(n)$ with little overhead in its circuit complexity and seed length, which is crucial for our applications. The idea is simple: we composite it with a hash with large shrinkage, short seed length, and relatively larger circuit complexity. In particular, we can use the following hash construction $H^{\mathrm{expw}}$ of Chen, Jin, and Williams [9]. (Here expw stands for expander walk.)
\begin{authorclaim}{Construction4 (Hash function $H^{\mathrm{expw}}$ from expander walk [9])}
Let $\ell=\ell(n)$ be the output length. $H^{\mathrm{expw}}_\ell=\{H_n\}_{\ge n}$ is constructed as follows.
\begin{itemize}
\item Let $\{S_n\}_{n\ge1}$ be a family of $0.1$-biased set from Theorem2.4 and $\{G_n\}_{n\ge1}$ be the family of strongly explicit expanders from Theorem2.10. Assume that $S_n=\{w_0,w_1,\ldots,w_{t-1}\}$ for $t=\widetilde O(n^2)$.
\item For each walk $v=(v_0,v_1,\ldots,v_{\ell-1})\in[t]^\ell$ of length $\ell$ on $G_t$, we define a hash function $h_v:\{0,1\}^n\to\{0,1\}^\ell$ as
\[h_v(x)\triangleq\langle w_{v_0},x\rangle\Vert\langle w_{v_1},x\rangle\Vert\cdots\Vert\langle w_{v_{\ell-1}},x\rangle,\]
where $\langle a,b\rangle$ denotes the inner product over $\mathbb F_2$.
\item $H_n$ is then defined to be the distribution generated as follows: samples a random walk $v$ of length $\ell$ on $G_t$ uniformly at random, then outputs $h_v$.
\end{itemize}
\end{authorclaim}
\begin{authorclaim}{Lemma3.5 (Properties of the hash construction $H^{\mathrm{expw}}$ from [9])}
Let $\ell=\ell(n)$ be the output length. The followings hold for $H^{\mathrm{expw}}_\ell=\{H_n\}_{n\ge1}$:
\begin{enumerate}
\item Every $H\in\operatorname{supp}(H_n)$ can implemented by a depth-1 $\mathsf{CC}^0[2]$ circuit of wire complexity $n\ell$ or a $B_2$ circuit of size $n\ell$ and depth $\log n+O(1)$.
\item $H^{\mathrm{expw}}_\ell$ is a linear $O(\log n+\ell)$-samplable $\ell$-output $\exp(-\Omega(\ell))$-almost universal hash function.
\end{enumerate}
\end{authorclaim}
\begin{proof}
The first item follows directly from the definition of $h_v$ from Construction4. In the following we establish the second item.
The linearity, seed length, and output length are straightforward to verify, thus we only analyze the collision probability. Since the hash function is linear, we only need to show that for any non-zero $x\in\mathbb F_2^n$, $\Pr_{h\leftarrow H_n}[h(x)=0]\le\exp(-\Omega(\ell))$. By Theorem2.4, we know that for any non-zero $x$, at least $0.4$ fraction of vectors $w\in S_n$ satisfies $\langle w,x\rangle=1$. Then by Theorem2.11, a random walk of length $\ell$ will find one of such $w$ with probability $1-\exp(-\Omega(\ell))$. This immediately implies that $h(x)\ne0$ with probability at most $\exp(-\Omega(\ell))$ for $h\leftarrow H_n$.
\end{proof}
Next we formally define the composition of two hash function families.
\begin{authorclaim}{Definition3.6 (Composition of hash families)}
Let $F=\{F_n\}_{n\ge1}$ and $G=\{G_n\}_{n\ge1}$ be two families of hash functions. The composition of $F$ and $G$ is defined as $F\circ G=\{(F\circ G)_n\}_{n\ge1}$, where $(F\circ G)_n$ is the following distribution: let $m=m(n)$ be the output length of $G$, we sample $f\leftarrow F_m$ and $g\leftarrow G_n$, and then output $f\circ g$.
\end{authorclaim}
The following proposition is crucial for the analysis of our final hash construction.
\begin{authorclaim}{Proposition3.7}
The composition $H'\circ H$ of an $\varepsilon_1$-almost universal hash function $H'$ and an $\varepsilon_2$-almost universal hash function $H$ is an $(\varepsilon_2(n)+\varepsilon_1(m(n)))$-almost universal hash function, where $m(n)$ is the output length of $H$. Moreover, $H'\circ H$ is linear if both of $H'$ and $H$ are linear.
\end{authorclaim}
The following is a simple corollary of Proposition3.7.
\begin{authorclaim}{Corollary3.8}
For any $t\ge2$, non-increasing $\varepsilon=\varepsilon(n)$, and output length parameter $\ell(n)\le n$, the $t$th order composition of an $\ell$-output $\varepsilon$-almost universal hash $H$ with itself, denoted by $H^{\circ t}$, is also a hash function with collision probability $t\cdot\varepsilon(\hat\ell)$, where $\hat\ell=\ell\circ\ell\circ\cdots\circ\ell(n)$ ($t-1$ times) is the input of the outer-most hash. Moreover, $H^{\circ t}$ is linear if $H$ is linear.
\end{authorclaim}
We are now ready to specify our final hash construction $H^{\mathrm{final}}$ with $\operatorname{polylog}(n)$ output length.
\begin{authorclaim}{Construction5 (Construction of $H^{\mathrm{final}}$ with $\operatorname{polylog}(n)$ output length)}
Let $\beta\in(0,2]$ be a parameter and $\omega(\log n)\le\ell(n)\le O(\log^2n/\log\log n)$ be a non-decreasing function. We define
\[H^{\mathrm{final}}_{\beta,\ell}\triangleq H^{\mathrm{expw}}_\ell\circ(H^{\mathrm{md}}_{\beta,\ell})^{\circ\lceil2/\beta\rceil}.\]
\end{authorclaim}
We analyze the properties of $H^{\mathrm{final}}$ constructed above by combining the composition proposition (Proposition3.7 and Corollary3.8) together with properties of $H^{\mathrm{expw}}$ (Lemma3.5) and $H^{\mathrm{md}}$ (Theorem3.4).

\begin{authorclaim}{Theorem3.9 (Properties of the final hash construction $H^{\mathrm{final}}$)}
Let $\beta\in(0,2]$ be a constant and $\omega(\log n)\le\ell(n)\le O(\log^2n/\log\log n)$ be a non-decreasing function. The followings hold for $H^{\mathrm{final}}_{\beta,\ell}=\{H_n\}_{n\ge1}$:
\begin{enumerate}
\item $H^{\mathrm{final}}_{\beta,\ell}$ is a linear $O(\ell\log^{\beta+1}n)$-samplable $\Theta(\hat\ell)$-output $\exp(-\Omega(\hat\ell))$-almost universal hash function, where $\hat\ell=\ell(\Theta(n/\log^{\beta\cdot\lceil2/\beta\rceil}n))$.
\item Every $H\in H_n$ can be computed by depth-$(1+\lceil2/\beta\rceil)$ $\mathsf{CC}^0[2]$ circuits of wire complexity $2n+O(n\hat\ell/\log^2n)$, or by $B_2$ circuits of size $2n+O(n\hat\ell/\log^2n)$ and depth $\log n+O(1)$.
\end{enumerate}
\end{authorclaim}

\begin{proof}
Let $H^{\mathrm{md}}=H^{\mathrm{md}}_{\beta,\ell}$ and $H^{\mathrm{expw}}=H^{\mathrm{expw}}_\ell$ be the hash functions described in Construction5. We also use $H_n^{\mathrm{md}}$ and $H_n^{\mathrm{expw}}$ to denote the distributions over $n$-input hash functions in $H^{\mathrm{md}}$ and $H^{\mathrm{expw}}$, respectively.
\paragraph{$H^{\mathrm{final}}$ is almost universal with $\operatorname{polylog}(n)$ output length.}
Let $\hat m$ be the output length of $(H^{\mathrm{md}})^{\circ\lceil2/\beta\rceil}$, and $\hat\ell\triangleq\ell(\hat m)=\ell(\Theta(n/\log^{\beta\cdot\lceil2/\beta\rceil}))$. By a simple induction and Theorem3.4 we know that for any $1\le d\le\lceil2/\beta\rceil$, $(H^{\mathrm{md}})^{\circ d}$ is a linear $O(\ell\log^{\beta+1}n)$-samplable $\Theta(n/\log^{\beta d}n)$-output $\exp(-\Omega(\hat\ell))$-almost universal hash function. In particular, when $d=\lceil2/\beta\rceil$, we know that $(H^{\mathrm{md}})^{\circ\lceil2/\beta\rceil}$ is a linear $O(\ell\log^{\beta+1}n)$-samplable $\Theta(n/\log^{\beta\cdot\lceil2/\beta\rceil}n)$-output $\exp(-\Omega(\hat\ell))$-almost universal hash function.
Now apply Proposition3.7 and Lemma3.5, we know that $H^{\mathrm{final}}=H^{\mathrm{expw}}\circ(H^{\mathrm{md}})^{\circ\lceil2/\beta\rceil}$ is a linear $O(\ell\log^{\beta+1}n)$-samplable $\Theta(\hat\ell)$-output $\exp(-\Omega(\hat\ell))$-almost universal hash function.
\paragraph{Complexity with $\mathsf{CC}^0[2]$ circuits.}
By Theorem3.4, every $H\in\operatorname{supp}(H_n^{\mathrm{md}})$ can be computed by a depth-1 $\mathsf{CC}^0[2]$ circuit with wire complexity $2n$. Also, by Lemma3.5, every $H\in\operatorname{supp}(H_n^{\mathrm{expw}})$ can be computed by a depth-1 $\mathsf{CC}^0[2]$ circuit with wire complexity $n\ell$. Since the output length of $(H^{\mathrm{md}})^{\circ d}$ is $\Theta(n/\log^{\beta d})$, the total wire complexity of a hash function from the support of $H^{\mathrm{expw}}\circ(H^{\mathrm{md}})^{\circ\lceil2/\beta\rceil}$ can be bounded by
\[2n+O\left(\sum_{d=1}^{\lceil2/\beta\rceil-1}\frac n{\log^{\beta d}n}\right)+\frac n{\log^{\beta\cdot\lceil2/\beta\rceil}n}\cdot\hat\ell=2n+O\left(\frac{n\hat\ell}{\log^2n}\right).\]
The depth complexity can be verified straightforwardly.
\paragraph{Complexity with $B_2$ circuits.}
We only analyze the depth since the analysis of the size complexity is similar to the analysis above. By Theorem3.4 and Lemma3.5, every $H\in\operatorname{supp}(H_n^{\mathrm{md}})$ can be computed by a $B_2$ circuit with depth $\beta\log\log n+O(1)$, and every $H\in\operatorname{supp}(H_n^{\mathrm{expw}})$ can be computed by a $B_2$ circuit with depth $\log n+O(1)$. Therefore the total depth of a hash function from the support of $H^{\mathrm{expw}}\circ(H^{\mathrm{md}})^{\circ\lceil2/\beta\rceil}$ is at most
\[\sum_{d=0}^{\lceil2/\beta\rceil-1}\left[\beta\log\log\left(\Theta\left(\frac n{\log^{\beta d}n}\right)\right)+O(1)\right]+\log\left(\Theta\left(\frac n{\log^{\beta\cdot\lceil2/\beta\rceil}n}\right)\right)=\log n+O(1).\]
\end{proof}
\subsection*{3.4 Explicitness of our construction}
Apart from the seed length, shrinkage and collision probability, the explicitness of the a hash function is also critical for many applications. That is, given a seed $v$ (which is usually much shorter than the input, say $|v|=\operatorname{polylog}(n)$), whether we can obtain information about the hash function corresponding to the seed $v$ efficiently. In this section we show that our hash constructions are indeed explicit in a strong sense, which is crucial for our application to hardness magnification.
\begin{authorclaim}{Definition3.10 (Locally explicit hash function)}
A family of $\operatorname{polylog}(n)$-samplable distributions over $m$-output linear functions (e.g., linear hash functions or 1-detectors) $H=\{H_n\}_{n\ge1}$ is said to be locally explicit if each input bit only influences $\operatorname{polylog}(n)$ output bits, and there exists an algorithm $A(n,v,i)$ running in $\operatorname{polylog}(n)$ time that returns the list of output bits influenced by the $i$th input bit in the hash function corresponding to the seed $v$.
\end{authorclaim}
Note that any linear function can be realized as $x\mapsto Mx$ for some transformation matrix $M$ over $\mathbb F_2$. Clearly, a locally explicit linear hash has a sparse transformation matrix, and we can efficiently list all 1-entries in any column of it. By performing a sparse matrix multiplication, we can immediately obtain the following proposition.
\begin{authorclaim}{Proposition3.11}
The composition of two locally explicit hash families is locally explicit.
\end{authorclaim}
Now we verify that our constructions of the randomness-efficient low-complexity hash functions are indeed locally explicit.
\begin{authorclaim}{Theorem3.12}
Let $k=k(n)=n^{o(1)}$ be the shrinkage parameter and $\log n\le r(n)\le n/2$ be the range. The 1-detector $L^{\mathrm{hg}}_{k,r}$ in Construction2 is locally explicit.
\end{authorclaim}
\begin{proof}
Recall that the 1-detector in Construction2 consists of two components: a depth-1 $\mathsf{CC}^0[2]$ circuit whose topology is determined by the high-girth graph from Theorem2.12, and a random permutation of the input bits according to the $k$-wise almost independent permutation from Theorem2.8. Given the seed $v$ of the 1-detector (i.e., the seed for the permutation) and an index $i$. According to Theorem2.8, we can obtain $\sigma_v(i)$ in $\operatorname{poly}(\log n,|v|)$ time, where $\sigma_v$ is the permutation corresponding to the seed $v$. To find all the output bits influenced by the $i$th input bit, we only need to find the two endpoints of the $\sigma_v(i)$-th edge in the high-girth graph, which can be done in $\operatorname{polylog}(n)$ time by Theorem2.12.
\end{proof}
\begin{authorclaim}{Theorem3.13}
Let $\beta>0$ be the shrinkage parameter and $\delta=\delta(n)\le O(\log^2n/\log\log n)$ be the error parameter. The hash family $H^{\mathrm{md}}_{\beta,\delta}$ in Construction3 is locally explicit.
\end{authorclaim}
\begin{proof}
Recall that the hash function $H^{\mathrm{md}}$ in Construction3 is obtained by the 1-detector $L^{\mathrm{hg}}$ in Construction2 and the reduction in Construction1. The output of the hash $H^{\mathrm{md}}_{\beta,\delta}$ consists of two parts: the output of the 1-detector $L^{\mathrm{hg}}$, and several randomly sampled bits from the input. Since $L^{\mathrm{hg}}$ is locally explicit by Theorem3.12, it is sufficient to show that the sampling procedure is locally explicit. More precisely, given a seed $v$ of the sampling procedure and an index $i$, we need to compute efficiently all the indices $j$ such that $w_j=i$ for the vector $(w_0,\ldots,w_{b-1})$ corresponding to the seed $v$. This can be done using the algorithm $A^{\mathrm{samp}}$ from Lemma3.1.
\end{proof}
\begin{authorclaim}{Theorem3.14}
Let $\ell=\ell(n)=\operatorname{polylog}(n)$. The hash family $H^{\mathrm{expw}}_\ell$ in Construction4 is locally explicit.
\end{authorclaim}
\begin{proof}
Let $S_n=\{w_0,\ldots,w_{t-1}\}$ with $t=\widetilde O(n^2)$ is a $0.1$-biased set from Theorem2.4 and $G_t$ is the expander in Theorem2.10. The hash function in Construction4 is defined as
\[h_v(x)\triangleq\langle w_{v_0},x\rangle\Vert\langle w_{v_1},x\rangle\Vert\cdots\Vert\langle w_{v_{\ell-1}},x\rangle,\]
where $v=(v_0,\ldots,v_{\ell-1})\in[t]^\ell$ is a random walk on $G_t$. Given the seed $v$ (i.e., the seed of a random walk), we can produce $v_0,v_1,\ldots,v_{\ell-1}$ in $\operatorname{polylog}(n,\ell)=\operatorname{polylog}(n)$ time by Theorem2.10. Then using the algorithm $A$ in Theorem2.4 we can check for each $j$ whether the $i$th bit of $w_{v_j}$ is 1, which indicates whether the $j$th output bit is influenced by the $i$th input bit. Enumerating all $j\in[\ell]$ gives a required $\operatorname{polylog}(n)$ time algorithm listing all the outputs influenced by a particular input bit.
\end{proof}
Using Theorem3.13, Theorem3.14, and Proposition3.11, we immediately obtain the explicitness of Construction5.
\begin{authorclaim}{Corollary3.15}
Let $\beta\in(0,2]$ be a parameter and $\ell=\ell(n)\le O(\log^2n/\log\log n)$ be a non-decreasing function. The hash family $H^{\mathrm{final}}_{\beta,\ell}$ in Construction5 is locally explicit.
\end{authorclaim}
\subsection*{3.5 Uniformity of our construction}
We have shown in Theorem3.9 that our hash function can be computed by extremely sparse $\mathsf{CC}^0[2]$ circuit or $B_2$ circuits with small size and depth simultaneously. It remains to clarify the uniformity of our hash construction, i.e., the complexity we need to construct the circuit given the input length and the seed.
The definitions of uniformity vary with respect to the complexity measures, which can be chosen according to the applications. Typical choices include space complexity (e.g., LOGSPACE-uniform, see Section6.2.1 of [3]), time complexity (e.g., polynomial-time uniform), parallel time (see, e.g., Section3 of [29]), and descriptive complexity (see, e.g., [20]). If we choose the time complexity, for instance, we can define P-uniformity as the existence of the polynomial-time algorithm that prints the circuit given $1^n$ and the seed $v$. Indeed, it is easy to check that the $2n+o(n)$ size $B_2$ and $\mathsf{CC}^0[2]$ circuits for Construction5 in Theorem3.9 is P-uniform.
Since our hash function can be evaluated in low-depth circuit models such as $\mathsf{NC}^1$ and $\mathsf{CC}^0[2]$, the parallel time to generate the circuit also seems to be crucial for its further applications. Therefore in this section we will discuss the POLYLOGTIME-uniformity of our hash functions, defined as follows.
\begin{authorclaim}{Definition3.16}
Let $C$ be a circuit class. A family of $\operatorname{polylog}(n)$-samplable distributions over $m$-output functions (e.g., hash functions or 1-detectors) $H=\{H_n\}_{n\ge1}$ is said to be computable with POLYLOGTIME-uniform $S(n)$-size (measured in size, depth, wire complexity, etc) $C$-circuits if each of the functions $h_v\in\operatorname{supp}(H_n)$ corresponding to the seed $v$ is computable by an $S(n)$-size $C$-circuit, supplemented with the following $\operatorname{polylog}(S(n))$-time algorithms:
\begin{itemize}
\item $\mathrm{SIZE}(n,v)$ returns the size of the $C$-circuit $C_v$ computing $h_v$;
\item $\mathrm{TYPE}(n,v,i)$ returns the type of the $i$th node in the circuit $C_v$, which indicates (1) whether it is a gate, an input variable or a constant; (2) whether it is an output node; and (3) the type (if it is a gate or a constant) or index (if it is an input variable and/or output node) of it;
\item $\mathrm{EDGE}(n,v,i,j)$ returns the $j$th input of the $i$th node if the $i$th node is a gate in $C_v$.
\item $\mathrm{OUT}(n,v,i)$ returns the number of the $i$th output node.
\end{itemize}
It is guaranteed that the nodes are numbered in topological orders, and in particular, the input variables are numbered from 1 to $n$.
\end{authorclaim}
We claim that all our results of circuit complexity for the constructions in the previous parts of this section are in fact POLYLOGTIME-uniform. It is straightforward to check the claim, so we only sketch the proof and left the details to the readers.
\begin{authorclaim}{Theorem3.17 (Uniformity of Theorem3.9)}
Let $\beta\in(0,2]$ be a parameter and $\ell=\ell(n)\le O(\log^2n/\log\log n)$ be a non-decreasing function. The hash function $H^{\mathrm{final}}_{\beta,\ell}$ in Construction5 can be computed either by POLYLOGTIME-uniform depth-$(1+\lceil2/\beta\rceil)$ $\mathsf{CC}^0[2]$ circuits of wire complexity $2n+O(n/\log\log n)$, or by POLYLOGTIME-uniform $B_2$ circuits of size $2n+O(n/\log\log n)$ and depth $\log n+O(\beta^{-1})$ simultaneously.
\end{authorclaim}
\paragraph{Sketch.}
Since the POLYLOGTIME-uniformity is close under composition, we only need to verify that Construction3 and Construction4 are POLYLOGTIME-uniform. The uniformity of Construction3 directly follows from the algorithm $B$ in Theorem2.12, the algorithm $A$ in Theorem2.8, and the algorithm $A$ in Lemma3.1. The uniformity of Construction4 follows from the algorithm $A$ in Theorem2.4 and the algorithm $A$ in Theorem2.10. Note that the circuit for Construction4 should be slightly modified to make $\mathrm{EDGE}(n,\cdot,\cdot)$ $\operatorname{polylog}(n)$-time computable. Take the $\mathsf{CC}^0[2]$ case for example: each of the output gates has fan-in exactly $n$ (instead of $|w_{v_i}|$ for the $i$th output bit, see Construction4 for the notation); the $j$th input wire of the $i$th output gate is connected to the $j$th input or a constant 0 according to the $j$th bit of $w_{v_i}$. \hfill$\square$

% Faithful typesetting transcription of AppendixA, formal pp23:36--23:37.
\section*{AppendixA: Strongly explicit high-girth graphs}
In this section, we will examine the construction of bipartite high-girth graphs of Lazebnik and Ustimenko [27], and verify that it is indeed strongly explicit, thereby proving Theorem2.12.
Let $q$ be an odd prime power. Let $P$ and $L$ be two infinite sequences of elements from $\mathbb F_q$ indexed as follows:
\begin{align*}
P&=p_1,p_{1,1},p_{1,2},p_{2,1},p_{2,2},p'_{2,2},p_{3,2},\ldots,p_{i,i+1},p_{i+1,i},p_{i+1,i+1},p'_{i+1,i+1},p_{i+1,i+2},\ldots,\\
L&=l_1,l_{1,1},l_{1,2},l_{2,1},l_{2,2},l'_{2,2},l_{3,2},\ldots,l_{i,i+1},l_{i+1,i},l_{i+1,i+1},l'_{i+1,i+1},l_{i+1,i+2},\ldots.
\end{align*}
The names $P$ and $L$ mean points and lines, respectively, due to certain geometric intuition of the construction. We say $P$ is incident with $L$ if they satisfy the following set of equations
\begin{equation}\tag{3}
E_1=\begin{cases}l_{1,1}-p_{1,1}=l_1p_1\\l_{1,2}-p_{1,2}=l_{1,1}p_1\\l_{2,1}-p_{2,1}=l_1p_{1,1},\end{cases}\qquad
E_i=\begin{cases}l_{i,i}-p_{i,i}=l_1p_{i-1,i}\\l'_{i,i}-p'_{i,i}=l_{i,i-1}p_1\\l_{i,i+1}-p_{i,i+1}=l_{i,i}p_1\\l_{i+1,i}-p_{i+1,i}=l_1p'_{i,i}\end{cases}\quad(\forall i\ge2).
\end{equation}
One can verify that for every point $P$, a line $L$ incident with it is uniquely determined by $l_1$, since all other coordinates of $L$ can be computed from the equations above iteratively. Similarly for every line $L$, a point $P$ incident with it is uniquely determined by $p_1$.
Let $P_k\in\mathbb F_q^k$ be the length-$k$ prefix of $P$, and $L_k\in\mathbb F_q^k$ be the length-$k$ prefix of $L$. We define $D(k,q)=(V_1,V_2,E\subseteq V_1\times V_2)$ to be the following bipartite graph:
\begin{itemize}
\item $V_1=V_2=\mathbb F_q^k$;
\item $(u,v)\in E\Longleftrightarrow\exists(P,L)$, $P$ is incident with $L$, $u=P_k$, and $v=L_k$.
\end{itemize}
Equivalently, $u$ and $v$ are connected if and only if they satisfy the first $k$ equations of $\bigcup_i E_i$ in (3). Note that $|V_1|=|V_2|=q^k$, $|E|=q^{k+1}$, and the graph is $q$-regular.
\begin{authorclaim}{TheoremA.1 (Lazebnik and Ustimenko [27])}
Let $k\ge3$ be an odd integer and $q$ be an odd prime power. Then the girth of $D(k,q)$ is at least $k+5$.
\end{authorclaim}
For our purpose, it is sufficient to consider $q=p^r$ for prime $p=O(1)$. To construct the graph explicitly, we need to evaluation field operations over $\mathbb F_q$ for a prime power $q$ in $\operatorname{polylog}(q)$ time, for which we need a representation of $\mathbb F_q$. Recall that $\mathbb F_q$ is isomorphic to $\mathbb F_p[x]/(Q(x))$ for any irreducible degree-$r$ polynomial $Q\in\mathbb F_p[x]$. So the explicit representation of $\mathbb F_q$ follows from the construction of irreducible polynomials by Shoup [39].
\begin{authorclaim}{TheoremA.2 (Shoup [39])}
There is a deterministic algorithm that constructs a degree-$n$ irreducible polynomial over $\mathbb F_p$ given an integer $n$ and a prime power $p$ in $\operatorname{poly}(n,p)$ time.
\end{authorclaim}
Now we check that $D(k,q)$ is strongly explicit, in the sense that given the index $i$ of an edge, we can obtain the indices $j_1$ and $j_2$ in $\operatorname{poly}(k,\log q)$ time such that the $i$th edge connects the $j_1$th vertex in $V_1$ and the $j_2$th vertex in $V_2$. Let $q=p^r$ for a prime $p=O(1)$. We number the vertices and edges as follows.
\begin{enumerate}
\item We identify the elements in the finite field $\mathbb F_q$ as length-$r$ vectors from $[p]^r$, and number all the elements in the lexicographic order.
\item We identify the vertices in both $V_1$ and $V_2$ (that are length $k$ sequences of elements in $\mathbb F_q$) as length-$kr$ vectors from $[p]^{kr}$, and number them in the lexicographic order.
\item For any $i\in[q^k]$ and $j\in[q]$, the $(iq+j)$th edge connects the $i$th vertex $P$ in $V_1$ and the unique vertex $L=(l_1,\ldots)$ in $V_2$ connected to $P$ such that $l_1$ is the $j$th element in $\mathbb F_q$.
\end{enumerate}
Under such a numbering scheme, given the index $iq+j$ of an edge with $i\in[q^k]$ and $j\in[q]$, we can easily determine its two endpoints in $\operatorname{poly}(k,\log q)$ time. Now we are ready to prove Theorem2.12.
\begin{authorclaim}{Remainder of Theorem2.12}
Let $r=r(n)=n^{o(1)}$ be a parameter. For every sufficiently large $n$, there exists an $m=\Theta(n/r)$ and a regular graph $G_{m,n}$ with $m$ vertices, $n$ edges, and girth $\Omega(\log n/\log r)$. Moreover, there exists a $\operatorname{polylog}(n)$-time algorithm $A(n,i)$ for $i\in[n]$ that outputs the indices of the two endpoints of the $i$th edge in $G_{m,n}$, and a $\operatorname{polylog}(n)$-time algorithm $B(n,i,j)$ for $i\in[m]$ that outputs the $j$th edge attaching to the $i$th vertex.
\end{authorclaim}
\begin{proof}
Let $q$ be the power of 3 in the interval $[r,3r)$. Applying TheoremA.2, we construct a fixed representation of $\mathbb F_q$ in $\operatorname{polylog}(q)$ time so that the field operations over $\mathbb F_q$ can be evaluated in $\operatorname{polylog}(q)$ time. Let $k\triangleq\lfloor\log_qn\rfloor-1$ (i.e., $k$ is the largest integer such that $q^{k+1}\le n$), $\ell\triangleq\lceil n/q^{k+1}\rceil$, and $m\triangleq2q^k\ell$. Note that
\begin{align*}
m&\le2q^k\left(\frac n{q^{k+1}}+1\right)=\frac{2n}q+\frac{2q^{k+1}}q\le O(n/r),\\
m&\ge2q^k\cdot\frac n{q^{k+1}}\ge\Omega(n/r),
\end{align*}
and therefore $m=\Theta(n/r)$. Now we construct a graph $G_{m,n}$ with $\ell$ connected components as follows. Each of the first $\ell-1$ connected components is a copy of $D(q,k)$ with $2q^k$ vertices and $q^{k+1}$ edges. The last connected component is a subgraph of $D(q,k)$ with $2q^k$ vertices but only the first $n-q^{k+1}(\ell-1)$ edges.
By TheoremA.1, it is easy to see that the girth of $G_{m,n}$ is at least $k=\Omega(\log n/\log q)=\Omega(\log n/\log r)$. According to our numbering scheme and related discussions, it is also easy to verify that given an edge index $i$, we can compute the two endpoints of the $i$th edge in $G_{m,n}$ in $\operatorname{polylog}(n)$ time. Furthermore, the $j$th edge attaching to the $i$th vertex can be easily computed since each of the $q$ edges attaching to the $i$th vertex is uniquely determined by $l_1\in\mathbb F_q$ or $p_1\in\mathbb F_q$ (see Equation3).
\end{proof}

\begin{thebibliography}{43}

\bibitem[1]{ref1}
Noga Alon, Oded Goldreich, Johan Håstad, and René Peralta. Simple constructions of almost k-wise independent random variables. In 31st Annual Symposium on Foundations of Computer Science, St. Louis, Missouri, USA, October 22-24, 1990, Volume II, pages 544–553. IEEE Computer Society, 1990. doi:10.1109/FSCS.1990.89575.

\bibitem[2]{ref2}
Benny Applebaum and Pavel Raykov. Fast pseudorandom functions based on expander graphs. In Theory of Cryptography - 14th International Conference, TCC 2016-B, Beijing, China, October 31 - November 3, 2016, Proceedings, Part I, volume 9985 of Lecture Notes in Computer Science, pages 27–56, 2016. doi:10.1007/978-3-662-53641-4\_2.

\bibitem[3]{ref3}
Sanjeev Arora and Boaz Barak. Computational Complexity - A Modern Approach. Cambridge University Press, 2009.

\bibitem[4]{ref4}
Abhishek Banerjee, Chris Peikert, and Alon Rosen. Pseudorandom functions and lattices. In Advances in Cryptology - EUROCRYPT 2012 - 31st Annual International Conference on the Theory and Applications of Cryptographic Techniques, Cambridge, UK, April 15-19, 2012. Proceedings, volume 7237 of Lecture Notes in Computer Science, pages 719–737. Springer, 2012. doi:10.1007/978-3-642-29011-4\_42.

\bibitem[5]{ref5}
Andrej Bogdanov and Alon Rosen. Pseudorandom functions: Three decades later. In Tutorials on the Foundations of Cryptography, pages 79–158. Springer International Publishing, 2017. doi:10.1007/978-3-319-57048-8\_3.

\bibitem[6]{ref6}
Larry Carter and Mark N. Wegman. Universal classes of hash functions. J. Comput. Syst. Sci., 18(2):143–154, 1979. doi:10.1016/0022-0000(79)90044-8.

\bibitem[7]{ref7}
L. Sunil Chandran. A high girth graph construction. SIAM J. Discret. Math., 16(3):366–370, 2003. doi:10.1137/S0895480101387893.

\bibitem[8]{ref8}
Lijie Chen, Shuichi Hirahara, Igor Carboni Oliveira, Ján Pich, Ninad Rajgopal, and Rahul Santhanam. Beyond natural proofs: Hardness magnification and locality. In 11th Innovations in Theoretical Computer Science Conference, ITCS 2020, January 12-14, 2020, Seattle, Washington, USA, volume 151 of LIPIcs, pages 70:1–70:48. Schloss Dagstuhl - Leibniz-Zentrum für Informatik, 2020. doi:10.4230/LIPIcs.ITCS.2020.70.

\bibitem[9]{ref9}
Lijie Chen, Ce Jin, and R. Ryan Williams. Hardness magnification for all sparse NP languages. In 60th IEEE Annual Symposium on Foundations of Computer Science, FOCS 2019, Baltimore, Maryland, USA, November 9-12, 2019, pages 1240–1255. IEEE Computer Society, 2019. doi:10.1109/FOCS.2019.00077.

\bibitem[10]{ref10}
Lijie Chen, Ce Jin, and R. Ryan Williams. Sharp threshold results for computational complexity. In Proccedings of the 52nd Annual ACM SIGACT Symposium on Theory of Computing, STOC 2020, Chicago, IL, USA, June 22-26, 2020, pages 1335–1348. ACM, 2020. doi:10.1145/3357713.3384283.

\bibitem[11]{ref11}
Ruiwen Chen and Valentine Kabanets. Correlation bounds and \#SAT algorithms for small linear-size circuits. Theor. Comput. Sci., 654:2–10, 2016. doi:10.1016/j.tcs.2016.05.005.

\bibitem[12]{ref12}
Evgeny Demenkov and Alexander S. Kulikov. An elementary proof of a 3n - o(n) lower bound on the circuit complexity of affine dispersers. In Mathematical Foundations of Computer Science 2011 - 36th International Symposium, MFCS 2011, Warsaw, Poland, August 22-26, 2011. Proceedings, volume 6907 of Lecture Notes in Computer Science, pages 256–265. Springer, 2011. doi:10.1007/978-3-642-22993-0\_25.

\bibitem[13]{ref13}
Zhiyuan Fan, Jiatu Li, and Tianqi Yang. The exact complexity of pseudorandom functions and the black-box natural proof barrier for bootstrapping results in computational complexity. Electron. Colloquium Comput. Complex., page 125, 2021. To appear in STOC 2022. URL: https://eccc.weizmann.ac.il/report/2021/125.

\bibitem[14]{ref14}
Magnus Gausdal Find, Alexander Golovnev, Edward A. Hirsch, and Alexander S. Kulikov. A better-than-3n lower bound for the circuit complexity of an explicit function. In IEEE 57th Annual Symposium on Foundations of Computer Science, FOCS 2016, 9-11 October 2016, Hyatt Regency, New Brunswick, New Jersey, USA, pages 89–98. IEEE Computer Society, 2016. doi:10.1109/FOCS.2016.19.

\bibitem[15]{ref15}
Lance Fortnow. A simple proof of Toda’s theorem. Theory Comput., 5(1):135–140, 2009. doi:10.4086/toc.2009.v005a007.

\bibitem[16]{ref16}
Oded Goldreich, Shafi Goldwasser, and Silvio Micali. How to construct random functions. J. ACM, 33(4):792–807, 1986. doi:10.1145/6490.6503.

\bibitem[17]{ref17}
Alexander Golovnev, Edward A. Hirsch, Alexander Knop, and Alexander S. Kulikov. On the limits of gate elimination. J. Comput. Syst. Sci., 96:107–119, 2018. doi:10.1016/j.jcss. 2018.04.005.

\bibitem[18]{ref18}
Johan Håstad. The shrinkage exponent of de Morgan formulas is 2. SIAM J. Comput., 27(1):48–64, 1998. doi:10.1137/S0097539794261556.

\bibitem[19]{ref19}
Johan Håstad, Russell Impagliazzo, Leonid A. Levin, and Michael Luby. A pseudorandom generator from any one-way function. SIAM J. Comput., 28(4):1364–1396, 1999. doi: 10.1137/S0097539793244708.

\bibitem[20]{ref20}
William Hesse, Eric Allender, and David A. Mix Barrington. Uniform constant-depth threshold circuits for division and iterated multiplication. J. Comput. Syst. Sci., 65(4):695–716, 2002. doi:10.1016/S0022-0000(02)00025-9.

\bibitem[21]{ref21}
Shuichi Hirahara and Rahul Santhanam. On the average-case complexity of MCSP and its variants. In 32nd Computational Complexity Conference, CCC 2017, July 6-9, 2017, Riga, Latvia, volume 79 of LIPIcs, pages 7:1–7:20. Schloss Dagstuhl - Leibniz-Zentrum für Informatik, 2017. doi:10.4230/LIPIcs.CCC.2017.7.

\bibitem[22]{ref22}
Russell Impagliazzo, Raghu Meka, and David Zuckerman. Pseudorandomness from shrinkage. J. ACM, 66(2):11:1–11:16, 2019. doi:10.1145/3230630.

\bibitem[23]{ref23}
Yuval Ishai, Eyal Kushilevitz, Rafail Ostrovsky, and Amit Sahai. Cryptography with constant computational overhead. In Proceedings of the 40th Annual ACM Symposium on Theory of Computing, Victoria, British Columbia, Canada, May 17-20, 2008, pages 433–442. ACM, 2008. doi:10.1145/1374376.1374438.

\bibitem[24]{ref24}
Kazuo Iwama and Hiroki Morizumi. An explicit lower bound of 5n - o(n) for boolean circuits. In Mathematical Foundations of Computer Science 2002, 27th International Symposium, MFCS 2002, Warsaw, Poland, August 26-30, 2002, Proceedings, volume 2420 of Lecture Notes in Computer Science, pages 353–364. Springer, 2002. doi:10.1007/3-540-45687-2\_29.

\bibitem[25]{ref25}
Eyal Kaplan, Moni Naor, and Omer Reingold. Derandomized constructions of k-wise (almost) independent permutations. Algorithmica, 55(1):113–133, 2009. doi:10.1007/ s00453-008-9267-y.

\bibitem[26]{ref26}
Oded Lachish and Ran Raz. Explicit lower bound of 4.5n - o(n) for boolena circuits. In Jeffrey Scott Vitter, Paul G. Spirakis, and Mihalis Yannakakis, editors, Proceedings on 33rd Annual ACM Symposium on Theory of Computing, July 6-8, 2001, Heraklion, Crete, Greece, pages 399–408. ACM, 2001. doi:10.1145/380752.380832.

\bibitem[27]{ref27}
Felix Lazebnik and Vasiliy A. Ustimenko. Explicit construction of graphs with an arbitrary large girth and of large size. Discret. Appl. Math., 60(1-3):275–284, 1995. doi:10.1016/ 0166-218X(94)00058-L.

\bibitem[28]{ref28}
Jiatu Li and Tianqi Yang. 3.1n - o(n) circuit lower bounds for explicit functions. Electron. Colloqu ium Comput. Complex., page 23, 2021. To appear in STOC 2022. URL: https: //eccc.weizmann.ac.il/report/2021/023.

\bibitem[29]{ref29}
Richard J. Lipton and Ryan Williams. Amplifying circuit lower bounds against polynomial time, with applications. Comput. Complex., 22(2):311–343, 2013. doi:10.1007/s00037-013-0069-5.

\bibitem[30]{ref30}
Dylan M. McKay, Cody D. Murray, and R. Ryan Williams. Weak lower bounds on resource- bounded compression imply strong separations of complexity classes. In Proceedings of the 51st Annual ACM SIGACT Symposium on Theory of Computing, STOC 2019, Phoenix, AZ, USA, June 23-26, 2019, pages 1215–1225. ACM, 2019. doi:10.1145/3313276.3316396.

\bibitem[31]{ref31}
Eric Miles and Emanuele Viola. Substitution-permutation networks, pseudorandom functions, and natural proofs. J. ACM, 62(6):46:1–46:29, 2015. doi:10.1145/2792978.

\bibitem[32]{ref32}
Ketan Mulmuley. On P vs. NP and geometric complexity theory: Dedicated to sri ramakrishna. J. ACM, 58(2):5:1–5:26, 2011. doi:10.1145/1944345.1944346.

\bibitem[33]{ref33}
Moni Naor and Omer Reingold. Number-theoretic constructions of efficient pseudo-random functions. J. ACM, 51(2):231–262, 2004. doi:10.1145/972639.972643.

\bibitem[34]{ref34}
Igor Carboni Oliveira, Ján Pich, and Rahul Santhanam. Hardness magnification near state-of- the-art lower bounds. Theory Comput., 17:1–38, 2021.

\bibitem[35]{ref35}
Igor Carboni Oliveira and Rahul Santhanam. Hardness magnification for natural problems. In 59th IEEE Annual Symposium on Foundations of Computer Science, FOCS 2018, pages 65–76, 2018. doi:10.1109/FOCS.2018.00016.

\bibitem[36]{ref36}
Anna Pagh and Rasmus Pagh. Uniform hashing in constant time and optimal space. SIAM J. Comput., 38(1):85–96, 2008. doi:10.1137/060658400.

\bibitem[37]{ref37}
Christos H. Papadimitriou and Stathis Zachos. Two remarks on the power of counting. In Theoretical Computer Science, 6th GI-Conference, Dortmund, Germany, January 5-7, 1983, Proceedings, volume 145 of Lecture Notes in Computer Science, pages 269–276. Springer, 1983. doi:10.1007/BFb0009651.

\bibitem[38]{ref38}
Alexander A. Razborov and Steven Rudich. Natural proofs. J. Comput. Syst. Sci., 55(1):24–35, 1997. doi:10.1006/jcss.1997.1494.

\bibitem[39]{ref39}
Victor Shoup. New algorithms for finding irreducible polynomials over finite fields. In 29th Annual Symposium on Foundations of Computer Science, White Plains, New York, USA, 24-26 October 1988, pages 283–290. IEEE Computer Society, 1988. doi:10.1109/SFCS.1988.21944.

\bibitem[40]{ref40}
Daniel A. Spielman. Linear-time encodable and decodable error-correcting codes. IEEE Trans. Inf. Theory, 42(6):1723–1731, 1996. doi:10.1109/18.556668.

\bibitem[41]{ref41}
Avishay Tal. Shrinkage of de Morgan formulae by spectral techniques. In 55th IEEE Annual Symposium on Foundations of Computer Science, FOCS 2014, Philadelphia, PA, USA, October 18-21, 2014, pages 551–560. IEEE Computer Society, 2014. doi:10.1109/FOCS.2014.65.

\bibitem[42]{ref42}
Leslie G. Valiant and Vijay V. Vazirani. NP is as easy as detecting unique solutions. Theor. Comput. Sci., 47(3):85–93, 1986. doi:10.1016/0304-3975(86)90135-0.

\bibitem[43]{ref43}
Emanuele Viola. The communication complexity of addition. Comb., 35(6):703–747, 2015. doi:10.1007/s00493-014-3078-3.
\end{thebibliography}
\end{document}
